\subsection{Search instructions, repair tasks, and retrieval of earlier objections}
\label{sec:complementary}
Three companion studies examine how counterexamples affect error search, formal proof repair, and proof revision; each study has its own input panel and evaluation protocol (Appendix~\ref{app:complementary-detail}).

\paragraph{Search instructions.}
We compared search instructions on 100 proofs flagged by the workflow as containing errors with concrete witnesses (an earlier panel distinct from the localization benchmark). Counterexample-first search instructs the model to look first for a refutable claim and a concrete counterexample. The comparison condition reviews the proof sequentially and requires a counterexample before stopping. A model grader, given the reference counterexample, counts an objection as \emph{confirmed} when it accepts the response's counterexample or self-contained refutation; a correct ``invalid'' verdict without such support does not count. Counterexample-first search used 31\%, 15\%, and 29\% fewer output tokens than sequential review with a required counterexample for GPT-5.6 Sol, Gemini 3.1 Pro, and Opus 5. The corresponding differences in confirmed-objection rates were $+2$, $-5$, and $+5$ percentage points, with all 95\% intervals including zero. Counterexample-first search used 6--11\% more output tokens than ordinary early stopping, which allows an objection without a counterexample. On a separate panel of 60 false mathematical statements presented without proofs, counterexample-first search used 26--52\% fewer output tokens than a proof attempt that must supply a counterexample when it rejects a statement (Table~\ref{tab:instruction-main}). Output-token counts include reported reasoning and do not measure end-to-end cost: ratios of recorded token totals depend on prompt caching and recorded tool executions (Section~\ref{app:primary-witness-accounting} and Appendix~\ref{app:cost-repro-procedures}).

\paragraph{Repair tasks.}
On 16 Lean repair tasks in which the correction to prove is given, we test whether supplying a checked counterexample helps a model complete the formal repair. Both conditions receive the source context, formal definitions, identified defect, repair template, and target theorem statements; the treatment condition adds the dossier's checked counterexample, and a fixed checker accepts submissions whose proofs establish the required statements (Section~\ref{app:repair-study}). On first attempts, adding the counterexample raised completed repairs from 13 to 15 for Sonnet~5 and from 1 to 2 for GPT-OSS 120B, while Opus~5.5 completed 15 of 16 in both conditions (Table~\ref{tab:repair_outcomes}). The study tests completing a supplied repair rather than discovering one from scratch. Because the task set is small and the effect varies by model, broader benefits remain to be tested.

\paragraph{Retrieval of earlier objections.}
Two studies tested giving a model access to earlier BugBank records while it worked on new problems. In a proof-revision study on 50 paired tasks per model under a shared stopping rule, a model grader judged 45 proofs correct with retrieval versus 47 without for Sol, 43 versus 42 for Gemini, and 44 versus 43 for Opus 5. Recorded token totals were higher with retrieval for Sol and Gemini and lower for Opus (Table~\ref{tab:proof50-results}, Section~\ref{app:proof-development}). In an exploratory open-problem study across 36 target problems, model audits accepted 12 resolutions with retrieval and 11 without, 11 of them shared by both conditions (Table~\ref{tab:historical_retrieval}, Section~\ref{app:historical-retrieval}). Retrieving stored objections thus produced mixed results; determining which prior objections transfer to a new argument, and when they justify the extra context, remains an open question.
